Researchers and practitioners use synthetic equity paths to test risk
models, evaluate portfolio strategies, and examine market scenarios
beyond those observed in a historical sample~\cite{assefaEtAl2020,
jordonEtAl2022}. A multi-asset simulator must represent both the
behavior of each asset and the relationships among assets. Factor
models provide a simple way to link many assets through a shared
market path, a time series of market returns~\cite{sharpe1963,
rossAPT1976, famaFrench1993}. The simplest of these, the single-index
model, regresses each asset's return on the market path. It separates
the return into an intercept, a market component equal to the market
path scaled by the fitted loading, and a zero-mean asset-specific
residual. A synthetic path for the asset can then be composed by
adding a simulated residual to the intercept and market component.
Suppose single-asset generators, models that simulate return paths
for one asset, have already been fitted to each asset's full returns
rather than to its regression residuals. Their simulated paths are then
an obvious candidate for that residual. Each such generator, however, already represents
its asset's total variance, including the part driven by the market. Adding its paths to the market component as
though they were residuals therefore counts market variance a second
time. The paths also generally have a nonzero mean, unlike residuals;
subtracting each path's mean removes that shift but leaves market
variance counted twice. Fitting a second generator to each asset's
regression residuals avoids the problem but repeats the model fitting for every asset. Reusing the existing generators would
avoid that second fit, provided their paths could be transformed to
supply only the residual contribution. The question is whether such a
transformation can control the mean and variance of the composed paths
while retaining the generator's heavy tails and volatility clustering.

The transformation is useful only if each generator reproduces the
behavior of its own asset. We illustrate the features a generator must
capture with the broad-market exchange-traded fund with ticker SPY,
which also supplied the market path in our multi-asset experiments.
Throughout the paper, including in the generators and the
regression, returns are excess growth rates. Each day's log price
change is multiplied by $252$ trading days to express it as an annual
rate, and a constant annual risk-free rate is subtracted. For SPY from $2014$--$2024$, these growth rates show heavy
tails and volatility clustering (Fig.~\ref{fig:empirical_motivation}).
Most growth rates lie close to zero, but extreme values occur much more
often than a Gaussian model predicts
(Fig.~\ref{fig:empirical_motivation}a,b), a widely documented feature
of financial returns~\cite{mandelbrot1963, cont2001}. Raw growth rates
show little correlation across days, while absolute growth rates remain
positively correlated across many lags
(Fig.~\ref{fig:empirical_motivation}c,d). Large and small absolute
growth rates therefore cluster in time, a property that synthetic
financial data often miss even when their return histogram appears
realistic~\cite{stengerEtAl2024}.

A hidden Markov model (HMM) offers an interpretable way to represent
both features: it draws growth rates from state-dependent distributions
and uses transitions between states to represent temporal
behavior~\cite{rabinerJuang1986, rydenTerasvirtaAsbrink1998}. In the
model used here, the states are ordered bins of growth rates, from the
most negative to the most positive. Each observed day's state is read
from its bin rather than inferred, and transition probabilities come
from counting day-to-day moves between bins. Each state draws growth
rates from a Student-$t$ distribution whose location and scale come
from the days in its bin; because that distribution is unbounded,
simulated values can fall outside the bin. The chain remembers only the current day's bin, not whether recent
days were volatile. Fitted to SPY, the model leaves
the extreme-return states at either end quickly, and the
autocorrelation of its absolute growth rates decays geometrically with
the lag, faster than in the data. A duration mechanism that
occasionally holds the model in the extreme-return states for several
consecutive days can slow that decay~\cite{bullaBulla2006}. Adding the
market component, however, changes the distribution and temporal
behavior of the composed returns, so an improvement from such a mechanism, whether added to SPY's
generator or to the asset generators, must be checked again after
composition.

% Main figure fig:empirical_motivation is collected after the references.

In this study, we developed a variance correction for combining
full-return generators, each fitted to one asset, with a shared market
path (Fig.~\ref{fig:preservation}; Table~\ref{tab:aggregate}). The
correction first centered each simulated full-return path by
subtracting its own mean over the simulated days. The centered path, used as the residual, therefore summed to zero over
those days, so the intercept and market component alone set the
composed mean. The same constraint removes asset-specific uncertainty
from each path's cumulative return over the simulated days. The correction then multiplied
the centered path by a factor below one. For each simulated path, the factor was chosen so that the market
component's variance plus the scaled path's variance equaled the path's variance before scaling, which represents the
asset's total variance. Because each asset path was simulated
independently of the market path, their sample covariance was small, and the composed variance approximately equaled
that total. We applied the correction to heavy-tailed hidden
Markov generators for the other $423$ assets in the $424$-asset
universe, first without the duration mechanism to isolate the
correction. The market path came from SPY: observed in the
$2014$--$2024$ training comparison, and simulated from SPY's fitted
generator in the $2025$ holdout so that no $2025$ data entered the
simulation. We compared the correction with naive composition, which
only centered the paths, and with Gaussian residuals whose variance
came from the regression. We also compared it with three methods that
fitted a separate generator to each asset's regression residuals: a
hidden Markov generator, a block bootstrap, and a generalized
autoregressive conditional heteroskedasticity model. This main comparison measured distributional fit, variance, and
Value-at-Risk thresholds. Volatility clustering of the composed paths, which the
duration mechanism targets, was measured in a separate comparison with
and without that mechanism.

The Kolmogorov--Smirnov pass rate measured distributional fit as the
share of simulated paths that a two-sample test at the $5\%$ level did
not distinguish from an asset's observed growth rates. In the training period, the correction passed the Kolmogorov--Smirnov
test far more often than naive composition and kept heavy tails,
whereas paths with Gaussian residuals had much lower kurtosis
(Table~\ref{tab:aggregate}).
The residual-fit hidden Markov generator had a slightly higher training
pass rate than the correction. With all $2014$--$2024$ fits held fixed,
we then compared the composed paths with the $416$ of the $423$ assets
that had complete $2025$ histories. Against naive composition, the
correction raised the mean pass rate across assets and brought the
median ratio of synthetic to observed variance from well above one to
close to one (Table~\ref{tab:oos_composition}). In this holdout, the
residual-fit hidden Markov generator had a slightly lower pass rate
than the correction, and every other comparator had a lower pass rate
than either. One-day Value-at-Risk checks at the $95\%$ and $99\%$
levels counted how often each asset's observed $2025$ losses exceeded
thresholds estimated from the simulated paths (Table~\ref{tab:var}).
For both the correction and the residual-fit hidden Markov generator,
exceedances were slightly more frequent than the nominal $5\%$ and
$1\%$ rates. Because the assets shared one market and holdout days were
not independent, we treated these rates as descriptive rather than as a
formal test of calibration.

We assessed the SPY hidden Markov generator in more detail, together
with a jump-duration variant that implements the duration mechanism
(Table~\ref{tab:model_comparison}). In this variant, a jump
is not an instantaneous price change but an episode. With a small
probability each day, the model entered an episode lasting a random
number of consecutive days. On each episode day, the state was drawn
from either the most negative or the most positive states, with the
side chosen at random each day. In the SPY training period, the variant had a smaller autocorrelation
error, the gap between simulated and observed autocorrelations of
absolute growth rates, than the model without jumps; in the $2025$
holdout, the two errors were similar. The main comparison excluded jumps, so the separate comparison tested
whether this training improvement carried over to the composed paths
(Table~\ref{tab:jump_ablation}). The main comparison used observed SPY in training, whereas this one
simulated the market path from the SPY generator in both periods, so
that jumps in that generator could affect the market component. It
contrasted no jumps with two configurations: jumps in the SPY generator
only, and jumps in that generator and every asset generator. Both
configurations used the jump parameters calibrated on SPY's
$2014$--$2024$ data, without recalibration for each asset, and each configuration was run
with both naive and corrected composition.

Both jump configurations reduced the autocorrelation error of the
composed paths in training and raised it in the $2025$ holdout
(Table~\ref{tab:jump_ablation}). Without jumps, the correction left
the autocorrelation error essentially where naive composition left it,
so the rescaling neither damaged nor improved clustering. With jumps,
the composed variance still matched the generator variance to within a
fraction of a percent on average. Across these experiments, the
correction removed the double-counted market variance, kept heavy
tails, and left clustering where naive composition left it. These
results support reuse as a practical option when full-return
generators already exist, with residual fitting as a competitive
alternative. Neither reuse nor residual fitting, as tested
here, modeled dependence among the residuals of different assets, which
portfolio evaluation also requires~\cite{embrechtsMcNeilStraumann2002,
boucheyNemtchinovWong2015}.
